\paragraph{Resultados preliminares y experimentales}

Los entrenamientos fueron agrupados por funciones de pérdida, en particular por una estructura de tipos dentro de estas para ayudar a su posterior documentación y evaluación:

\begin{itemize}
    \item \textbf{Pérdidas principales}: $L1$, $L2$, $L1\_freq$, $\log L1$, $\log L2$, $\log\_mag$, $\log\_compressed\_l2$, $LPSA$, $mask\_L1$.
    \item \textbf{Pérdidas experimentales}: deep\_feature, deep\_feature\_emd.
    \item \textbf{Pérdidas avanzadas}: l\_mrs, lpsa\_phase.
\end{itemize}

También existe la división que ya se explicó anteriormente entre dos y cuatro fuentes. El caso de las dos fuentes representa una separación más sencilla y cercana a escenarios de voz y acompañamiento.
El caso de cuatro fuentes introduce mayor complejidad, ya que el modelo debe repartir la mezcla entre varias fuentes instrumentales con solapamiento espectral.

A tener en cuenta sobre la comparativa entre funciones de pérdida del proyecto:

\begin{itemize}
    \item En las primeras versiones del entrenamiento se detectaron problemas de formato, inferencia y evaluación. El proyecto terminó habiéndose realizado tres tandas de entrenamientos diferentes corrigiendo estos errores, de las cuales la última es la definitiva.
    \item Las pérdidas experimentales basadas en características profundas, tienen mayor coste y mayor incertidumbre, porque dependen de su extractor auxiliar.
    \item Las pérdidas avanzadas con fase o reconstrucción temporal son más complejas y, por tanto, conviene añadir control sobre ellas.
\end{itemize}